\documentclass[border=8pt,tikz]{standalone}
\usepackage{fontspec}
\setmainfont{Times New Roman}
\setsansfont{Arial}
\usepackage{polyglossia}
\setdefaultlanguage{russian}
\setotherlanguage{english}
\usepackage{tikz}
\usetikzlibrary{positioning,arrows.meta,calc,decorations.pathreplacing}

\definecolor{ink}{HTML}{1A1A1A}
\definecolor{mid}{HTML}{555555}
\definecolor{light}{HTML}{999999}
\definecolor{g1}{HTML}{D4D4D4}
\definecolor{g2}{HTML}{A3A3A3}
\definecolor{g3}{HTML}{737373}
\definecolor{g4}{HTML}{525252}
\definecolor{g5}{HTML}{E5E5E5}

\tikzset{
  every node/.style={font=\sffamily},
  arr/.style={-{Stealth[length=6pt,width=5pt]}, line width=.8pt, color=ink},
  note/.style={font=\sffamily\scriptsize, text=mid},
  bignote/.style={font=\sffamily\small\bfseries, text=ink},
}

\begin{document}
\begin{tikzpicture}

% ===== LEFT: Raw argmax =====
\node[bignote, anchor=south west] at (0, 3.2) {(а) Сырые argmax-предсказания};

% Noisy bar — many tiny random segments
\foreach \x/\shade/\w in {
  0/g5/.22, .22/g2/.35, .57/g1/.15, .72/g3/.28,
  1/g5/.18, 1.18/g2/.4, 1.58/g1/.12, 1.7/g3/.22,
  1.92/g5/.15, 2.07/g2/.08, 2.15/g1/.3, 2.45/g3/.15,
  2.6/g2/.2, 2.8/g5/.1, 2.9/g1/.25, 3.15/g3/.2,
  3.35/g2/.15, 3.5/g5/.25, 3.75/g1/.1, 3.85/g3/.3,
  4.15/g2/.2, 4.35/g5/.15, 4.5/g1/.25, 4.75/g3/.25
}{
  \fill[\shade] (\x, 2.4) rectangle (\x+\w, 3);
}
\draw[ink, line width=.4pt] (0, 2.4) rectangle (5, 3);

\node[note, anchor=north west] at (0, 2.3) {28 сегментов, Boundary F1 = 0.31};

% Problem annotations
\node[note, anchor=north west, text width=5cm] at (0, 1.8) {
  — Интро в середине трека\\
  — Осцилляции Припев$\leftrightarrow$Бридж\\
  — Сегменты длиной до 1\,с
};

% ===== ARROW =====
\draw[arr, line width=1.2pt] (5.5, 2.7) -- node[above, font=\sffamily\footnotesize\bfseries] {Витерби}
                                             node[below, font=\sffamily\tiny, text=light] {$\lambda{=}1.0$ + позиц. приоры}
                              (7.3, 2.7);

% ===== RIGHT: After Viterbi =====
\node[bignote, anchor=south west] at (7.5, 3.2) {(б) После Витерби + приоров};

% Clean segments
\fill[g5] (7.5, 2.4) rectangle (8.2, 3);    % Intro
\fill[g2] (8.2, 2.4) rectangle (9.5, 3);    % Verse
\fill[g3] (9.5, 2.4) rectangle (10.5, 3);   % Chorus
\fill[g4] (10.5, 2.4) rectangle (11.2, 3);  % Bridge
\fill[g3] (11.2, 2.4) rectangle (12.2, 3);  % Chorus
\fill[g5] (12.2, 2.4) rectangle (13, 3);    % Outro

\draw[ink, line width=.4pt] (7.5, 2.4) rectangle (13, 3);

% Labels inside
\node[font=\sffamily\tiny, text=ink] at (7.85, 2.7) {Интро};
\node[font=\sffamily\tiny, text=white] at (8.85, 2.7) {Куплет};
\node[font=\sffamily\tiny, text=white] at (10, 2.7) {Припев};
\node[font=\sffamily\tiny, text=white] at (10.85, 2.7) {Бр.};
\node[font=\sffamily\tiny, text=white] at (11.7, 2.7) {Припев};
\node[font=\sffamily\tiny, text=ink] at (12.6, 2.7) {Аутро};

\node[note, anchor=north west] at (7.5, 2.3) {7 сегментов, Boundary F1 = 0.45};

\node[note, anchor=north west, text width=5cm] at (7.5, 1.8) {
  — Интро только в начале, Аутро в конце\\
  — Согласованный Куплет$\to$Припев\\
  — Мин. длительность сегмента 6\,с
};

% ===== BOTTOM: delta =====
\node[draw=ink, line width=.8pt, rounded corners=3pt,
      minimum height=.7cm, text width=4cm, align=center,
      font=\sffamily\bfseries, text=ink,
      anchor=north] at (6.5, .6) {$\Delta$ Boundary F1: $+0.14$ ($+$45\%)};

\end{tikzpicture}
\end{document}
